\section{A second result: Pell doubling and one product square}
\label{sec:doubling-square}

The second construction has ten strictly positive integer unknowns.
It supplies two cost--degree tradeoffs: an 84-operation polynomial of
exact degree 436, and an 85-operation polynomial of exact degree 328.
Both retain the complete four-integer coefficient tuple of the inspected
89-operation, degree-288 parent.  In particular, this is a measured
reduction in arithmetic length, with a degree increase; it is not a
smaller-variable or global-minimality claim.

An incoming eleven-integer construction already uses a single product
square, enforcing coprimality by multiplying a strong Pell ordinate by
the full base and main radicand.  We credit that prior construction.
The change here is to \emph{double a Pell pair before eliminating its
roots}.  The resulting polynomial identity makes the auxiliary modulus
a multiple of the entire combined radicand.  Coprimality then follows
without that extra saturation.  The classical doubling identity itself
is not new; the claim is this application and the two explicit equations.

\subsection{The complete polynomial front}

The unknowns are precisely
\[
 J,F,\alpha,i,T,w,Y,n,\eta,\zeta\in\mathbb Z_{>0}.
\]
The four \emph{fixed} coefficient integers $\beta,K_0,M_C,M_F$ are
the whole directly indexed BFS943 tuple retained from the pinned source.
All four literal compressed byte strings, their bit lengths and hashes,
and their symbolic specification are included with the circuits.  They
satisfy
\[
 \beta\ge31\text{ odd},\qquad K_0>0,\qquad
 0<M_C<\beta<M_F<2\beta.
\]
Every symbol below is a polynomial abbreviation, with no further
existential quantifier:
\begin{align*}
 q&=\beta J+1,&h&=q-1,&C&=q-F-\alpha,\\
 X&=wq,&\mu&=q^2+1,&k&=\eta+\zeta,\\
 c&=kY+\eta,&A&=\mu Y(X+1)+1,&\Delta&=A^2-1,\\
 H&=4A-5,&V_0&=X\mu^2Y^2,&E&=1+V_0(V_0+1)k^2,\\
 R&=[h(q-F)+\alpha](q^2-1)+(M_C+qM_F)J,\\
 e&=k-R-1,&B&=(K_0+w)C-F,&U&=ic^2,\\
 a&=\Delta U^2E,&F_0&=1+a,&G&=EF_0,\\
 F_{\rm sq}&=1+4aF_0,&Q&=\Delta(F_{\rm sq}-1),\\
 b&=c(T+\alpha)-R,&V&=F_{\rm sq}b-c,\\
 K_+&=X(3c-2X)+2,&P&=q^3h,&M&=P\mu^2H.
\end{align*}
Choose $L=V^2$ for the 84-operation equation, or $L=bV$ for the
85-operation equation.  Put
\begin{align*}
 Z&=P\mu^2\{H(n-LB)+K_+\}
       +H\{\mu^2(hY-q^3B)-Pe\},\\
 R_G&=M^2G.
\end{align*}
The single equation is
\begin{equation}
 \boxed{Q(Z^2-R_GV^2)+(R_G-Z^2)=0.}
 \label{eq:doubling-final}
\end{equation}
The circuit files are \texttt{equation\_84.json} and
\texttt{equation\_85.json}.  Their operators are exclusively binary
addition, subtraction and multiplication.

\subsection{Why the product square recovers every required root}

For this proof only, write $\delta=Q-1$ and $D_y=QV^2-1$.
Equation~\eqref{eq:doubling-final} is the identical polynomial
\begin{equation}
 \delta Z^2-M^2GD_y=0.
 \label{eq:doubling-product}
\end{equation}
The early positive-input bounds give $A>4$, $\Delta\ge3$, $E>1$,
$c\ge3$, $U\ge c^2$, $F_{\rm sq}>c+1$, $Q>1$ and $M>0$.
Initially $b$ may have either sign.  If $b=0$, then $V=-c$;
otherwise $|V|\ge F_{\rm sq}-c>1$.  Therefore $D_y>0$ and
$\delta>0$ before any sign assertion about $B$ or $C$.

The essential polynomial identities are
\begin{equation}
 F_{\rm sq}=(1+2a)^2,\qquad
 Q=4\Delta^2U^2G.
 \label{eq:doubling-gcd}
\end{equation}
Thus $\delta\equiv D_y\equiv-1\pmod G$, giving
$\gcd(G,\delta D_y)=1$.  At a zero, set $z=Z/M\in\mathbb Q$.
The equation $\delta z^2=GD_y$ forces every prime valuation in $G$
to be even.  This argument is valid even if $M$ shares prime factors
with $G$, because a rational square still has even valuations.
Consequently $G=g^2$ for a positive integer $g$.

Since $F_0=1+\Delta U^2E$ is coprime to $E$, the factors of
$G=EF_0$ are separately squares:
\[
 E=\tau^2,\qquad F_0=f_0^2,\qquad g=\tau f_0.
\]
The positive rational $y=|z|/g$ satisfies
\begin{equation}
 QV^2-(Q-1)y^2=1.
 \label{eq:doubling-adjacent}
\end{equation}
It is integral.  To see this, write $y=r/m$ in lowest terms and put
$v=|V|$.  Clearing denominators gives $m^2\mid Q-1$.  For $v>1$,
$v<y<Qv/(Q-1)$, and the norm-preserving transformation
\[
 v'=(2Q-1)v-2(Q-1)y,\qquad
 y'=(2Q-1)y-2Qv
\]
keeps $v'$ integral and the reduced denominator $m$ unchanged, while
$-v<v'<v$.  Reflecting signs and descending on $|v|$ reaches $v=1$,
where $y=1$.  Hence $m=1$, and $z=\pm gy$ is integral.

Define the strong pair in the proof by
\[
 f=1+2a,\qquad I=2Ug=2U\tau f_0.
\]
Then
\[
 f^2-\Delta I^2=1,\qquad Q=(\Delta I)^2,\qquad c^2\mid I,
 \qquad V=f^2\{c(T+\alpha)-R\}-c.
\]
It is the Pell double of the seed pair $(f_0,U\tau)$.  These are
exactly the strong and auxiliary interfaces needed by the retained
source theorem.

\subsection{The sign argument is established before decoding}

The inherited all-index folding theorem gives $c=\psi_A(p)$ for odd
$p\ge1$.  The first norm gives
$k=2\psi_{2V_0+1}(n_0)$.  Its first index is impossible, because it
would put $c=2Y+1$ strictly between consecutive Pell ordinates $1$
and $2A$.  Hence
\[
 k\ge8V_0+4,\qquad c>8q^5,\qquad c>X,\qquad K_+>0.
\]
The unchanged mask bound $R<q^4+(\alpha+2)q^2$ then gives $b>c\ge3$,
without assuming the sign of $B$.

Put
\[
 \kappa=K_+/H>0,\quad s=Y/q^3>0,\quad
 \rho=e/\mu^2,\quad t=B/h,\quad W=1+hL.
\]
The complete cleared numerator gives the rational identity
\begin{equation}
 z=n+\kappa+s-Wt-\rho.
 \label{eq:doubling-carrier}
\end{equation}
Moreover $b-e=c(T+\alpha)-k+1>0$, so $\rho<b/\mu^2$.

For $L=V^2$, the elementary inequality
\[
 a^2-g^2
   =E\{(\Delta U^2)E(\Delta U^2-1)-1\}>0
\]
gives $g<a$.  Since $F_{\rm sq}-1=4aF_0>a$,
$V>(F_{\rm sq}-1)b>3g$ and $V>b$.
Equation~\eqref{eq:doubling-adjacent} gives $y<3V/2$, hence
$gy<V^2/2$.  If $B\le-1$, then $Wt<-V^2$ and
\[
 n<\frac{V^2}{2}-V^2+\frac{V}{\mu^2}<0,
\]
a contradiction.

For $L=bV$, a sharper bound follows from the fifth auxiliary index.
Since $p$ is odd and $c\ge3$, one has $p\ge3$, so
$c\ge\psi_A(3)=4\Delta+3$.  It follows that $V>4Q$.
Writing $Q=S^2$, the positive auxiliary Pell trace has odd index;
indices one and three give $V=1$ and $V=4Q-3$ and are excluded.
Therefore
\[
 V\ge16Q^2-20Q+5>8Q^2,
\]
and
\[
 b>\frac{8\Delta^2(F_{\rm sq}-1)^2}{F_{\rm sq}}
   >4\Delta^3I^2>4g.
\]
Thus $gy<3bV/8<L/2$.  A hypothetical $B\le-1$ yields
$n<L/2-L+b/\mu^2<0$, again impossible.  Both versions therefore
give $B\ge0$, and its literal definition forces $C>0$.

At this point the inherited rank, parity, denominator and whole-product
projection arguments apply in their original order.  They require
$c^2\mid I$, the same fixed-minus residues, an integral carrier $z$,
and $W\equiv1\pmod h$.  All have been established.  The strong index
need not be odd, so doubling causes no parity obstruction.  The source
reconstruction consequently gives a finite syntactic ZFC refutation.

\subsection{Completeness and the two precise tradeoffs}

A finite refutation supplies the canonical outer, first and main data.
Choose a seed Pell period whose ordinate $I_0$ is divisible by
$c^2\tau$; existence follows from the finite order of the determinant-one
Pell matrix modulo $c^2\tau$.  Put $i=I_0/(c^2\tau)$ and double the
seed pair:
\[
 f=2f_0^2-1,\qquad I=2f_0I_0.
\]
The unchanged auxiliary-completeness theorem permits this strong pair
and gives a fresh positive $T$ and $y>V$.  It places no parity
restriction on the strong index or $I/c^2$.
The integer choice
\[
 n=gy+Wt+\rho-\kappa-s
\]
is positive: the canonical bounds give $\tau>Y\ge s$ and
$f_0>Xc>\kappa$, hence $g>\kappa+s$; also $Wt,\rho\ge0$.
Its carrier is $z=gy$, so the polynomial vanishes.  This completes
both existential directions.  With the explicitly retained ordinary
native calibration, ZFC proves absence of positive zeros equivalent
to $\operatorname{Con}(\mathrm{ZFC})$.

The 84-operation graph has 50 products, 19 sums and 15 differences;
the 85-operation graph adds the single product $bV$.  The exact degrees
are certified by global upper bounds and nonzero leading forms.
Writing $q_0=\beta J$ and $k_0=\eta+\zeta$, the degree-328 leading
part is
\[
 -2^{10}q_0^{156}w^{34}Y^{78}i^{14}k_0^{44}(q_0-F)^2,
\]
while the degree-436 leading part is
\[
 2^{14}q_0^{200}w^{46}Y^{104}i^{20}k_0^{60}
 (q_0-F)^4(q_0-F-\alpha)^2.
\]
Both are nonzero at the actual $\beta>0$.
The standalone standard-library verifier checks every literal gate,
all four coefficient bindings, the exact identities, degree upper
bounds and nonzero finite-field specializations using the actual tuple.
These finite checks supplement the ordinary proof; they do not replace
the retained unbounded calibration or constitute a new Lean proof object.
